\chapter{SAT and combinatorial constructions}
\label{ch:sat}
A molecular search does not evaluate a formula by understanding its symbols.
It changes a population so that the molecules left in one tube encode
assignments satisfying that formula. The difficult step is not writing
\emph{parallel} beside an arrow. It is establishing what each molecule means,
which copies can cross that arrow, and why a detected product is evidence
about the original problem.

Our running instance has three Boolean variables and three clauses:
\begin{equation}
 F=(x_1\lor x_2)\land(\neg x_1\lor x_3)
       \land(\neg x_2\lor\neg x_3).
 \label{eq:formula}
\end{equation}
Before reading further, predict the surviving assignments. Would a molecule
satisfying both literals in the first clause need to be copied twice? Would
an empty final tube prove that the formula is unsatisfiable?

Prerequisites are Chapter 11's injective encodings and Chapter 12's multiset
semantics. We will construct an exact clause-selection algorithm, prove its
invariant, implement it, and then remove the assumptions that make its
conclusion reliable. No sequence library or laboratory protocol is implied.

\section{From a formula to an addressable molecule}
\index{CNF|see{conjunctive normal form}}\index{assignment encoding}
A \Term{literal} is a variable or its negation. A \Term{clause} is a
disjunction of literals; a formula in \Term{conjunctive normal form} (CNF) is
a conjunction of clauses. An assignment is a vector
$a=(a_1,\ldots,a_n)\in\{0,1\}^n$. It satisfies a positive literal $x_i$
when $a_i=1$, and a negative literal $\neg x_i$ when $a_i=0$.
The empty disjunction is false; the empty conjunction is true. Those are
useful boundary cases, not syntax to discard.

Choose an addressable domain $D_{i,b}$ for variable $i$ and value $b$.
An ideal assignment strand contains
\begin{equation}
 E(a)=D_{1,a_1}D_{2,a_2}\cdots D_{n,a_n}.
\end{equation}
The ordered, fixed-width domain convention makes decoding unambiguous.
It does not establish biochemical orthogonality. A probe for $D_{2,1}$
must distinguish that domain from every other domain and from junctions,
secondary structures and contaminants. Chapter 11's sequence screens are
necessary checks, not proof that an extraction has zero error.

\Plate{01-encoding}{An assignment has stable variable-value addresses.
The probe pairs antiparallel to one address; the recognition event does not
change the represented assignment. Domain blocks are symbolic regions,
not individual bases or experimentally validated oligonucleotides.}{fig:satencoding}

There are $2^n$ possible assignments, although the address vocabulary uses
only $2n$ domains. Small vocabulary is not small inventory. Constructing
every possible concatenation requires coverage of the combinatorial space;
it does not follow from possessing one copy of every domain.
Lipton's 1995 paper proposed DNA computations for SAT \cite{lipton}.
The explicit algorithm below is an original finite-multiset teaching
construction in that tradition, not a reconstruction of an uninspected
experimental protocol.

\section{Implement an OR without inventing copies}
\index{clause partition}\index{multiset conservation}
Let $P(a)$ be the nonnegative integer number of molecules encoding $a$.
For a literal $\ell$, ideal separation creates two disjoint inventories:
\begin{equation}
 P_\ell^+(a)=P(a)\mathbf1[a\models\ell],\qquad
 P_\ell^-(a)=P(a)\mathbf1[a\not\models\ell].
 \label{eq:split}
\end{equation}
Thus $P_\ell^++P_\ell^-=P$ pointwise. Separation consumes the input tube
as an available inventory: drawing two output tubes must not leave a
third usable copy of the original.

For clause $C=(\ell_1\lor\cdots\lor\ell_k)$, start with residual $R_0=P$.
At step $j$, separate the residual, not the original:
\begin{align}
 A_j(a)&=R_{j-1}(a)\mathbf1[a\models\ell_j],\\
 R_j(a)&=R_{j-1}(a)\mathbf1[a\not\models\ell_j],\\
 Q(a)&=\sum_{j=1}^{k} A_j(a).
 \label{eq:clause}
\end{align}
An assignment takes its first matching branch. Once captured, it cannot
enter another branch. The final residual contains the assignments that
make every literal false.

\Plate{02-partition}{Clause selection is a disjoint decomposition of one
inventory. Arrows carry molecules, not copies of a software variable.
For $x_1\lor x_2$, assignment 110 is retained on the first branch and never
retested in the residual.}{fig:partition}

The common incorrect implementation extracts each literal from the full
original pool and adds the results. Seven copies of 110 then become
fourteen: both literals are true. It computes the right \emph{set} of
assignments, but the wrong physical multiset. If later readout weights
abundance, this error changes both the resource estimate and the evidence.
Actually dividing the original tube into aliquots avoids creating mass,
but no longer guarantees that each assignment reaches a matching branch.
That is a different randomized algorithm requiring its own completeness
analysis.

The companion function implements Equation~\ref{eq:clause}; signed
one-based integers encode literals. The public solver validates literal
addresses and copy counts before invoking it.
\Listing{clause_partition}

\section{Prove the invariant, then trace every assignment}
\index{completeness!relative to inventory}
Let $P_0$ be the initial population and $C_1,\ldots,C_m$ the clauses.
After $r$ clauses, the intended invariant is
\begin{equation}
 P_r(a)=P_0(a)\prod_{j=1}^{r}\mathbf1[a\models C_j].
 \label{eq:invariant}
\end{equation}
For one clause, an assignment either has a first true literal, and exactly
one $A_j$ receives all its copies, or has none, and remains in $R_k$.
Therefore $Q(a)=P(a)\mathbf1[a\models C]$. Applying this identity to
$P_{r-1}$ proves Equation~\ref{eq:invariant} by induction.

This proves \Term{soundness}: any surviving assignment satisfies every
processed clause. It proves completeness only \emph{relative to the initial
inventory}: a satisfying assignment with $P_0(a)=0$ cannot appear.
If each assignment begins with seven copies, our example progresses as
follows. The table is generated by the reference, not manually transcribed.

\begin{center}\input{\Generated/trace-table.tex}\end{center}

\Plate{03-cube}{Each vertex is one assignment, with $x_1x_2x_3$ as its
label. The nested survivor sets are an intersection of clause predicates;
they are not a time-dependent molecular conformation. Only 010 and 101
survive all three tests.}{fig:satcube}

A second implementation enumerates integer bit patterns and evaluates the
formula directly. It shares neither the sequential extraction algorithm
nor its candidate enumeration. Agreement with this oracle is necessary,
but not sufficient: both could still embody the same misunderstood
specification. The complete truth table makes the specification inspectable.

\begin{center}\input{\Generated/truth-table.tex}\end{center}
\Listing{solve}

An empty clause kills every assignment immediately. A repeated literal
does not change a clause because its second branch sees no new matches.
A clause containing both $x_i$ and $\neg x_i$ preserves the entire pool.
With zero variables and no clauses, the empty assignment is a valid witness.
These cases distinguish a Boolean solver from code that merely works for
one attractive demonstration.

\section{Where the exponential cost reappears}
\index{sampling coverage}\index{SAT!resource cost}
If every assignment has $r$ copies, the initial population contains
$r2^n$ molecules. If each encoded strand has $n$ domains of length $d$
bases, the assignment backbones alone contain
\begin{equation}
 B_{\mathrm{backbone}}=r2^n nd
\end{equation}
bases. Here $B$ is an inventory count, not a concentration or a mass.
Probe strands, linkers, purification losses and control samples are omitted.
The formula is a lower accounting boundary for this representation, not
a universal lower bound for all DNA algorithms.

For $m$ clauses with at most $k$ literals each, the ideal schedule uses at
most $mk$ separations and $m(k-1)$ binary merges (assuming $k\geq1$).
This is a polynomial number of bulk operations on an exponentially large
population. It is not a polynomial-resource solution to SAT. Moreover, bulk
operations do not have constant elapsed time independent of volume,
transport, capacity and the required error rate.

Uniform independent sampling offers a different tradeoff. With $s$
satisfying assignments among $2^n$, one sampled molecule is a witness with
probability $s/2^n$. If a witness survives the entire assay with probability
$q$, its probability of yielding evidence is $u=sq/2^n$. For $M$ samples:
\begin{align}
 \Pr(\text{no observed witness})&=(1-u)^M,\\
 \Pr(\text{at least one})&=1-(1-u)^M.
 \label{eq:coverage}
\end{align}
These equations assume independent sampling, uniform assignment abundance,
and independent survival with the same $q$ for every witness. Sequence bias
can violate all three assumptions.

\Listing{witness_probability}
\Plate{04-coverage}{Computed detection probability for a unique witness
among $2^{10}$ assignments, with perfect subsequent survival.
One sample per possible assignment \emph{on average} is not complete
coverage: random samples may repeat the same assignment.}{fig:coverage}

For a unique witness and $q=1$, taking $M=2^n$ gives approximately
$1-e^{-1}\approx0.632$, not certainty. To make miss probability at most
$\delta\in(0,1)$ when $0<u<1$, require
\begin{equation}
 M\geq\frac{\log\delta}{\log(1-u)}.
\end{equation}
Round upward to an integer. When $u=0$, no finite $M$ helps; when $u=1$,
one sample suffices. Computation uses \texttt{log1p} and \texttt{expm1}
to avoid subtracting nearly equal floating-point numbers.

\section{A positive witness and a negative result are asymmetric}
\index{witness verification}
The final readout should identify an assignment, not merely report bulk
fluorescence. Decode a strand into $a$ and evaluate $F(a)$ electronically.
This inexpensive certificate check can reject an invalid candidate caused
by cross-hybridization or contamination. It cannot recover a true witness
lost earlier. A verified witness proves satisfiability even if most of the
molecular computation malfunctioned.

\Plate{05-certificate}{A decoded assignment is checked against the
original formula, not merely against the last probe. A valid witness is
positive evidence. No signal remains compatible with either logical
unsatisfiability or failure to sample, retain or detect a witness.}{fig:certificate}

A negative claim requires stronger controls: positive recovery controls,
blank and contamination controls, a sequence-specific loss estimate,
a readout detection model, and evidence about initial coverage. Even then,
the conclusion may be a statistical upper bound on abundance rather than
a deductive UNSAT certificate. Do not convert an assay's detection threshold
into a theorem about all assignments.

Clause order is irrelevant to exact Equation~\ref{eq:invariant}, but can
matter physically. Early selective clauses reduce later load; early
damaging extractions may preferentially lose rare witnesses. Chapter 12's
capacity counterexample already shows why commuting Boolean predicates
need not yield commuting laboratory operators.

\section{Current research: change the architecture, not the accounting}
A September 2026 Scaffolded DNA Computer study reports computation through
engineered equilibrium preference, including finite-state programs and
25-bit addition \cite{sdc}. Its main text distinguishes short-scaffold
fast anneals from larger computations taking hours. It is not an
experimental demonstration of the exhaustive CNF tube algorithm above.

The useful research question is architectural: can a problem's structure
be compiled into local constraints so that one need not prepare every
candidate explicitly? A successful answer must still account for
preparation, degeneracy of incorrect states, equilibration and readout.
An energetically favored correct state does not, by itself, establish a
short path to it. Nor does a finite addition demonstration establish
efficient solution of arbitrary NP-complete instances.

\section{Exercises with worked answers}
\begin{enumerate}
\item \textbf{Predict.} Solve Equation~\ref{eq:formula} by cases on $x_1$.
\textit{Answer:} If $x_1=0$, clause 1 forces $x_2=1$, and clause 3 forces
$x_3=0$. If $x_1=1$, clause 2 forces $x_3=1$, and clause 3 forces $x_2=0$.
The witnesses are 010 and 101.
\item \textbf{Inventory.} A pool has seven copies of 110 and five of 000.
Apply $x_1\lor x_2$. \textit{Answer:} The first branch captures seven;
the second captures zero; five remain rejected. Naive independent extraction
from the original reports fourteen retained copies and violates conservation.
\item \textbf{Proof.} Why are the $A_j$ pairwise disjoint?
\textit{Answer:} Membership in $A_j$ requires all earlier literals false.
Membership in any earlier $A_i$ requires one of those literals true.
Both cannot hold for the same assignment.
\item \textbf{Boundary.} Compare no clauses with one empty clause.
\textit{Answer:} No clauses impose no restriction; an empty clause is false
for every assignment. With $n=0$, these return one empty witness and no
witness, respectively.
\item \textbf{Counterexample.} Can the ideal algorithm prove global
completeness from an arbitrary nonempty initial pool?
\textit{Answer:} No. A pool containing only 000 is nonempty, but has neither
witness of our formula. The proof is relative to represented assignments.
\item \textbf{Resource derivation.} With $n=20$, $r=10$, $d=15$, how many
assignment-backbone bases are required by the explicit inventory model?
\textit{Answer:} $10\cdot2^{20}\cdot20\cdot15=3{,}145{,}728{,}000$.
This omits all supporting reagents and says nothing about viable sequences.
\item \textbf{Sampling.} Explain the 0.632 limit.
\textit{Answer:} For $M=2^n$ and one witness,
$(1-2^{-n})^{2^n}\to e^{-1}$ is the miss probability.
A deterministic one-copy-per-assignment construction is not IID sampling.
\item \textbf{Loss.} A witness has ten independent retention steps of 0.9.
What changes in Equation~\ref{eq:coverage}?
\textit{Answer:} Substitute $q=0.9^{10}\approx0.348678$;
the effective success probability per draw decreases by that factor.
Correlated or sequence-dependent losses need a different model.
\item \textbf{Implementation.} Test all ordered pairs of two-literal clauses
over two variables, allowing repetition. \textit{Answer:} There are
$4^2=16$ ordered clauses and $16^2=256$ formulas. Compare the solver's
survivor keys with the integer oracle and check retained multiplicities.
The companion tests perform this exhaustive small check.
\item \textbf{Verification.} Does one recovered 101 certify that all
separations worked correctly? \textit{Answer:} No. It certifies SAT after
independent checking, not assay fidelity, yield or complete enumeration.
\item \textbf{Experimental design.} Design a no-signal investigation.
\textit{Rubric:} Specify a known witness spike-in, a negative sequence control,
a process blank, independent sequence identification, replicate loss estimates
and a coverage assumption. State what each failure would invalidate.
\item \textbf{Research critique.} Evaluate a claimed molecular SAT speedup.
\textit{Rubric:} Require instance distribution, full preparation and readout
costs, success probability, molecule budget, classical baseline, repeated
trials and failure cases. A count of separation rounds alone earns no
performance conclusion.
\end{enumerate}

\section{What we have built}
\textbf{Reproduce the chapter.} From the repository root, run
\texttt{python3 drafts/ch13/build.py --assets-only}, then
\texttt{python3 -m pytest tests/test\_ch13\_reference.py}.
The first command regenerates \texttt{results.json}, tables, plots and literal
code listings; the tests compare the implementation with independent
oracles and explicit failure cases. Run the builder without
\texttt{--assets-only} to compile this review PDF.

We now have an exact CNF selector whose inventory and logical invariants
can be checked independently. Its strength is a transparent correspondence
between predicates and physical-copy partitions. Its limitation is that
coverage, retention and measurement are additional obligations.

\textbf{Working vocabulary.} A literal tests one variable; a clause accepts
any of its literals; CNF requires all clauses; a witness is a satisfying
assignment; completeness relative to a pool is weaker than coverage of
all solutions. These distinctions recur when the stored assignment can
itself change. Chapter 14 replaces fixed sequence choices with a mutable
occupancy register and asks what it means to write, clear and reuse memory.
